\section{Track key identity within each map}
Use the decoded unsigned key to identify a field. Hexadecimal \texttt{00} and
\texttt{18 00} are different byte sequences for the same key zero. Comparing their raw
head bytes would miss a duplicate. With keys restricted to zero through seven, a
per-map bit mask can record which values have appeared.

The root encounters key three first, so its mask becomes hexadecimal \texttt{08}.
The optional map starts with a separate zero mask. Its keys one and zero produce mask
\texttt{03}; those keys do not conflict with keys one and zero that appear later in
the root. After returning, the root still owns mask \texttt{08}, then accepts keys two,
one and zero to reach \texttt{0f}.

Before publication, require the root's three required bits to be set. A mask of
\texttt{0b}, for example, includes keys zero, one and three but lacks the readings
field at key two. Reject duplicates rather than silently choosing a first or last value.
Reset the mask for each nested map and for every new top-level invocation.
